\section{Eerste uitwerking zwaartekrachtcompensatie via Euler-rotatiematrices}
\label{app:zwaartekracht}

Tijdens de ontwikkeling van de zwaartekrachtcompensatie is eerst een 
uitwerking gemaakt op basis van volledige Euler rotatiematrices. In deze 
aanpak wordt de volledige versnellingsvector via twee opeenvolgende 
rotaties teruggedraaid naar het horizontale referentievlak. Hierbij worden zowel 
de pitch als de rollhoek meegenomen. Hoewel deze methode gebaseerd is op de correcte wiskunde, bleek tijdens het testen dat er alsnog foutieve waardes ontstonden. De uiteindelijke implementatie gebruikt dan ook alleen de pitchhoek, zoals beschreven in hoofdstuk~\ref{sec:implementatie}. Deze bijlage bevat 
de volledige wiskundige uitwerking van de oorspronkelijke methode ter 
referentie.

Om de voorwaartse rijversnelling te bepalen, wordt de versnellingsvector 
$\mathbf{a}^b = [a_x, a_y, a_z]^T$ uit het gekantelde assenstelsel van 
de tracker teruggedraaid naar het horizontale referentievlak. Dit wordt 
gedaan met behulp van Euler-rotatiematrices. Zoals beschreven in de 
applicatienotitie \textit{Tilt Sensing Using a Three-Axis Accelerometer} 
van NXP Semiconductors~\cite{nxp2013}, wordt de rotatie om de $x$-as 
(Roll, $\phi_r$) en de $y$-as (Pitch, $\theta_p$) gedefinieerd door de 
volgende twee matrices:

\begin{equation}
R_x(\phi_r) = \begin{bmatrix} 
1 & 0 & 0 \\ 
0 & \cos\phi_r & \sin\phi_r \\ 
0 & -\sin\phi_r & \cos\phi_r 
\end{bmatrix}, \quad
R_y(\theta_p) = \begin{bmatrix} 
\cos\theta_p & 0 & -\sin\theta_p \\ 
0 & 1 & 0 \\ 
\sin\theta_p & 0 & \cos\theta_p 
\end{bmatrix}
\end{equation}

De totale transformatie naar het horizontale vlak wordt verkregen door 
de ruwe sensordata te vermenigvuldigen met deze matrices:

\begin{equation}
    \mathbf{a}_{\text{horizontaal}} = R_y(\theta_p) \cdot R_x(\phi_r) 
    \cdot \mathbf{a}^b
    \label{eq:horizontaal_bijlage}
\end{equation}

De berekening voor de voorwaartse as ($x$-component) wordt in twee 
stappen uiteengezet.

\paragraph{Stap 1: De roll rotatie ($R_x$)}

Eerst wordt de rollmatrix vermenigvuldigd met de sensordata. Omdat dit 
een draaiing om de $x$-as betreft, blijft de $x$-waarde zelf 
onveranderd, maar mixen de $y$- en $z$-waarden met elkaar. Dit geeft 
een tussenvector $\mathbf{v}'$:

\begin{equation}
    \mathbf{v}' = R_x(\phi_r) \cdot \begin{bmatrix} a_x \\ a_y \\ a_z 
    \end{bmatrix} = \begin{bmatrix} a_x \\ a_y \cos\phi_r + a_z 
    \sin\phi_r \\ -a_y \sin\phi_r + a_z \cos\phi_r \end{bmatrix}
    \label{eq:roll_matrix_bijlage}
\end{equation}

\paragraph{Stap 2: De pitch rotatie ($R_y$)}

Vervolgens wordt de pitchmatrix toegepast op de tussenvector 
$\mathbf{v}'$. Omdat alleen de voorwaartse as relevant is, wordt 
uitsluitend de bovenste rij van $R_y$ gebruikt, namelijk 
$[\cos\theta_p,\ 0,\ {-\sin\theta_p}]$:

\begin{equation}
    a_{\text{long}} = (\cos\theta_p \cdot v'_x) + (0 \cdot v'_y) - 
    (\sin\theta_p \cdot v'_z)
    \label{eq:pitch_matrix_bijlage}
\end{equation}

Na substitutie van de bekende waarden uit $\mathbf{v}'$ volgt:

\begin{equation}
    a_{\text{long}} = a_x \cos\theta_p - \sin\theta_p 
    (-a_y \sin\phi_r + a_z \cos\phi_r)
\end{equation}

\begin{equation}
    a_{\text{long}} = a_x \cdot \cos\theta_p + a_y \cdot \sin\phi_r 
    \cdot \sin\theta_p - a_z \cdot \cos\phi_r \cdot \sin\theta_p
    \label{eq:accel_correctie_bijlage}
\end{equation}

De termen met $a_y$ en $a_z$ compenseren het deel van de zwaartekracht 
dat door de hellingshoeken in de meting van de $x$-as is gelekt. In 
code zag deze uitwerking er als volgt uit:

\begin{lstlisting}[language=C]
float pitch_rad = ((float)current_attitude.pitch_x100 / 100.0f)
                  * ((float)M_PI / 180.0f);
float roll_rad  = ((float)current_attitude.roll_x100  / 100.0f)
                  * ((float)M_PI / 180.0f);

float flat_x = (raw_ax * cosf(pitch_rad)) +
               (raw_ay * sinf(roll_rad) * sinf(pitch_rad)) -
               (raw_az * cosf(roll_rad) * sinf(pitch_rad));
\end{lstlisting}

Tijdens het testen bleek dat het weglaten van de rollterm een 
verwaarloosbaar effect had op de nauwkeurigheid, omdat de tracker zo 
is gemonteerd dat de rollhoek klein blijft tijdens normaal gebruik. 
Op basis hiervan is de implementatie vereenvoudigd naar de methode 
beschreven in sectie~\ref{sec:implementatie}.